\documentclass[12pt, titlepage]{article}

\usepackage[T1]{fontenc}
\usepackage[utf8]{inputenc}
\usepackage[polish]{babel}

\usepackage{float}
\usepackage{graphicx}
\usepackage{tikz}
\usetikzlibrary{calc, arrows.meta, decorations.pathmorphing}
\usepackage{amsmath}
\usepackage{pdfpages}
\usepackage{caption}
\usepackage[backend=biber, style=numeric]{biblatex}

\captionsetup[table]{skip=6pt}
\addbibresource{refs.bib}

% typesetting
\usepackage{url}
\usepackage[hidelinks]{hyperref}
\usepackage{diagbox}
\usepackage{upgreek}
\usepackage{chemmacros}
\usepackage[version=4]{mhchem}
\usepackage{indentfirst}
\usepackage{geometry}
\usepackage{setspace}
\usepackage{makecell}
\usepackage{multirow}
\usepackage{pdfcomment}
\usepackage[normalem]{ulem}

\newgeometry{tmargin=2.5cm, bmargin=2.5cm, lmargin=3.5cm, rmargin=2.5cm}

\onehalfspacing
\tolerance 1414
\hbadness 1414
\emergencystretch 1.5em
\hfuzz 0.3pt
\widowpenalty=10000
\vfuzz \hfuzz
\raggedbottom
\hyphenpenalty=100000
% end packages

% title
\title{Sensory Laboratorium L5\\
Termopara}
\author{Wojciech Stawicki, Adam Burakowski, Krystian Cichewicz}
\date{Czerwiec 2026}

\begin{document}
\maketitle
\setcounter{page}{2}

\section{Wstęp teoretyczny}

\section{Cel ćwiczenia}

\section{Metodologia}

\section{Analiza danych}

\begin{figure}[H]
	\centering
	\includegraphics[width=\textwidth]{"plots/calibration.pdf"}
	\caption{Krzywa kalibracyjna}
	\label{fig:cal}
\end{figure}

\begin{figure}[H]
	\centering
	\includegraphics[width=\textwidth]{"plots/heating.pdf"}
	\caption{Grzanie}
	\label{fig:heat}
\end{figure}

\begin{figure}[H]
	\centering
	\includegraphics[width=\textwidth]{"plots/cooling.pdf"}
	\caption{Chłodzenie}
	\label{fig:cool}
\end{figure}

\section{Wnioski}

\end{document}
